% !TeX root = ../main.tex

\chapter{基于时空注意力机制-自适应图网络的交通流量预测方法}
\section{引言}

本文在上一章中介绍了一种基于时空注意力机制和图注意力网络的交通流预测的方法。在上一章的模型中, 研究使用了图网络和空间注意力机制来构建道路中的空间相关性, 使用了时间注意力机制来学习序列的时间相关性。虽然上一章的模型在交通流预测任务中, 通过融合注意力机制和优化网络结构取得了不错的成绩, 但是 GraphWavenet \parencite{2019_Wu_Graph_WaveNet_A_deep_learning_framework_for_traffic_forecasting} 和 AGCRN 这一类网络在 METR-LA 和 PEMS-BAY 数据集上的实验结果表现, 这类模型所设计的动态变化邻接矩阵在预测效果中有着较大的正向作用。 另外, 上一章节的部分消融模型 （仅使用时空注意力机制, 不使用图卷积的模型） 在 SZ 数据集上的精度较原始模型略高, 也说明了根据预先定义的邻接矩阵进行图卷积的运算, 在某些数据集的效率并不高。

目前主流交通流预测任务的都是基于图神经网络的时空网络模型进行预测。虽然各种基于图网络进行改进的模型已经在预测任务上取得了不错的效果, 但是这些模型仍存在以下问题:

\begin{enumerate}[(1)]
  \item 上述模型图卷积的部分, 从网络设计中可以看出所有节点参数的共享模式, 共享参数尽管能够降低参数量级并提取跨节点共性模式，但对于某些特定任务而言并非最优解。当前主流方法在精细化捕捉数据源特异性模式方面表现欠佳，这种现象源于参数共享策略的固有缺陷。交通流量序列本身具有多样化特征，不同数据源属性差异导致其模式呈现显著异质性甚至相互矛盾。由此可见，仅依赖共享模式提取难以满足预测精度要求。必须为每个独立节点配置专属参数集，节点特定模式的学习任务才能得以实现。

  \item 当前研究中，采用图卷积网络（GCN）的技术方案普遍存在着前置性约束条件。预先构建相互连接的图结构，该结构需借助相似度或距离度量的计算方式得以实现，其目的在于空间相关性特征的捕捉。这一过程不仅需要投入大量专业领域知识储备，更表现出对图结构质量的高度敏感性特征。由此可见，通过此类方法所生成的图结构往往呈现出直观性强但完整性不足的显著特点，且难以实现与预测任务需求间的直接适配性;\parencite{2021_Han_Dynamic_and_multi_faceted_spatio_temporal_deep_learning_for_traffic_speed_forecasting}。另外, 所有的图卷积操作都需要定义卷积核, 现有的卷积核有两种主要的定义方式: 首先是基于图拉普拉斯矩阵的卷积核, 该方法基于图的拉普拉斯矩阵来计算卷积核, 通常涉及到对图的特征值分解。具体来说, 图卷积操作通过拉普拉斯矩阵的特征向量来进行谱卷积。这种方式能够捕捉到图的全局结构信息, 并且为每个节点计算特征的加权和。然而, 这种方法的计算成本非常高, 尤其是在图特别大的情况下, 这将导致显著的资源消耗和时间延迟。为了解决基于拉普拉斯矩阵方法的计算效率问题, 另一种方法是设计不依赖于拉普拉斯矩阵特征向量的卷积核。在这种方法中, 卷积核需要通过巧妙设计来捕捉图的局部结构信息, 而不需要显式地计算图的谱分解。这样的设计通常依赖于图的邻接矩阵、图的局部拓扑结构、邻居节点的权重等信息。尽管这种方法在计算上更加高效, 但卷积核的设计通常需要丰富的领域知识和经验, 需要对图的结构有深入的理解, 并且可能涉及复杂的启发式搜索或者基于学习的优化方法。
\end{enumerate}

因此本文针对上述问题提出了一种基于注意力机制和自适应图网络的交通流预测方法。算法基于时空注意力机制-图神经网络的交通流预测研究, 提出了基于自适应图的图神经网络交通流预测模型, 以学习交通道路在空间中随着时间变化的连通属性。


\section{自适应图的设计思路}

为了解决图网络目前存在的上述两个问题, 本节将提出模型改进的设计思路。 自适应模块已经逐渐开始被应用于交通流量预测领域, FOGS 采用学习驱动的图结构, 使得模型可以动态调整不同道路节点间的关系, 提高预测的自适应性 \parencite{2022_Rao_FOGS_First_order_gradient_supervision_with_learning_based_graph_for_traffic_flow_forecasting}。\textcite{2021_Han_Dynamic_and_multi_faceted_spatio_temporal_deep_learning_for_traffic_speed_forecasting_2} 也在模型框架中引入可学习的图结构, 利用历史数据自适应调整道路间的连接权重, 从而更精准地刻画道路间的动态影响。本文的自适应模块主要包含两部分: 一个节点自适应参数学习 (NAPL) 模块, 用于捕获节点特定的模式; 以及一个数据自适应图生成(DAGG)模块, 以自动推断不同流量序列之间的相互依赖关系。

上一章的图的卷积运算思路, 主要基于下式 (展示为一层的图卷积运算):
\[
H = \sigma \left( {\theta XW}\right), 
\]
其中 \(\theta\) 是根据图拉普拉斯矩阵以及图滤波器理论设计的卷积核, \({X} \in  {\mathbb{R}}^{{N} * {C}}\) , 其中 \({N}\) 为节点数量, \({C}\) 为输入通道数量。 \({H} \in  {\mathbb{R}}^{{N} * {F}}\) , 其中 \({F}\) 为输出通道数量。 \({W} \in  {\mathbb{R}}^{{C} * {F}}\) , 是可学习的参数。从一个节点 (如节点 \({i}\) ) 的角度来看, \({{GCN}}\) 操作可以看作是将节点 \({{Xi}} \in  {\mathbb{R}}^{1 * {C}}\) 的特征转换为 \({{Hi}} \in  {\mathbb{R}}^{1 * {F}}\) , 所有节点之间共享参数 \({W}\)。来自两个相邻节点的流量流也可能在某些特定的属性或是某些特定时期呈现不同的模式。因此, 仅捕获所有节点之间的共享模式对于准确的流量预测是不够的。

但是, 为每个节点分配参数会导致 \({W} \in  {\mathbb{R}}^{{N} * {C} * {F}}\) , 参数量过大将导致计算成本的上升, 也会导致过拟合问题, 特别是当 \({N}\) 很大时。为解决该问题，本文创新性地引入节点自适应参数学习模块化组件。通过该模块与上一章模型有效整合，参数学习过程不再直接针对高维权重矩阵 \({W} \in {\mathbb{R}}^{{N} * {C} * {F}}\) 进行优化操作，转而采用两个低秩矩阵的协同学习机制: 1)节点嵌入矩阵 \({{E}}_{{g}} \in  {\mathbb{R}}^{{N} * {d}}\) , 其中 \({d}\) 为嵌入维数, \({d} \ll  {N}\) ; 通常可以选择 2)一个权重池 \({{W}}_{{g}} \in  {\mathbb{R}}^{{d} * {C} * {F}}\)。 然后, 可以通过 \({W} = {{E}}_{{g}} \cdot  {{W}}_{{g}}\) 生成 \({W}\)。从一个节点 (例如, 节点 \({i}\) ) 的角度来看, 该过程根据节点嵌入的 \({{E}}_{{{g}}_{i}}\) 从一个大型共享权重池 \({{W}}_{{g}}\) 中提取参数 \({{Wi}}\) , 这可以解释为从所有流量系列中发现的一组候选模式中学习出特定于节点的模式。最后, 节点自适应参数学习增强的 GCN 可以公式化为:
\[
H = \sigma \left( {{\theta X}{{E}}_{{g}} \cdot  {{W}}_{{g}}}\right).
\]

现存研究领域中，基于图卷积网络（GCN）的交通流预测模型存在着显著缺陷性。这类模型必须依赖预先定义的邻接矩阵A才能实施图的卷积运算操作。传统研究通常采用距离函数或相似性度量的计算方式提前构建图结构，但该方法难以捕捉节点间潜在的隐含关联性。针对这一局限性，本项研究工作创新性地设计出数据自适应的图生成机制，该模块能够从原始数据中自主推导出隐藏的相互依赖关系。具体实现过程中，首先进行的是所有节点的随机初始化处理，建立可学习的节点嵌入字典 \({{E}}_{{A}} \in  {\mathbb{R}}^{{N} * {{d}}_{{e}}}\) ，其中\({{E}}_{{A}}\) 的每一行向量对应着特定节点的嵌入表示，\({{d}}_{{e}}\) 则表征着节点嵌入的维度数量。与基于节点相似度的传统构图方法相仿，通过将\({{E}}_{{A}}\) 和 \({{E}}_{{A}}{}^{{T}}\) 进行乘积运算，可以推导出各节点对之间的空间依赖性关系。在此过程中，采用softmax函数对生成的自适应矩阵实施归一化处理。区别于传统方法需要生成邻接矩阵A并计算拉普拉斯矩阵的操作流程，本方案直接生成图卷积核 \(\theta\)。这种改进有效避免了迭代训练过程中冗余计算的产生。随着训练过程的推进，\({{E}}_{{A}}\) 将实现自动更新优化，从而学习到不同交通流量序列间潜在的依赖模式，最终形成适用于图卷积运算的自适应矩阵。相较于人工设计的图卷积核函数，这种自适应图生成机制展现出更简洁的结构特性，降低计算复杂度的同时提升了模型的表征能力。最后, 自适应图生成模块与节点自适应参数学习增强的 GCN 可以表述为:
\[
H = \sigma \left( {\operatorname{softmax}\left( {\operatorname{RELU}\left( {{{E}}_{{A}} \cdot  {{E}}_{{A}}{}^{{T}}}\right) }\right) X{{E}}_{{g}} \cdot  {{W}}_{{g}}}\right).
\]
因此, 将原网络层中的 GCN 替换为自适应图生成模块与节点自适应参数学习增强的 GCN, 生成时空注意力机制-自适应图网络。

\section{实验结果分析}

第 \ref{chap:based_gnn_attention} 章使用了 METR-LA 数据集和 SZ-taxi 数据集, 本章依旧在这 2 个数据集上进行实验。节点自适应参数学习中, \({{E}}_{{g}}\) 的参数 \({d}\) 设置为 \(2, {{E}}_{{A}}\) 的嵌入维度设置为 64, 即每个节点将学习一个 64 维度的节点表示。其他实验参数设置如第 \ref{chap:based_gnn_attention} 章介绍, 也不做赘述。本部分研究将只对比时空注意力机制-自适应图网络与上一章节的时空注意力机制-自适应图网络模型, 研究依旧采取相同的评估标准。

\begin{table}[htbp]
  \caption{自适应图网络与原模型对比}
  \centering
  \begin{tblr}{
    hline{1, Z} = {1pt}, 
    hline{2}, 
    width = \textwidth, 
    colspec = {X[1, c]X[3, c]X[1, c]X[1, c]}, 
    cell{2}{1} = {r=2}{m}, 
    cell{4}{1} = {r=2}{m}, 
  }
  && RMSE & MAE & MAPE \\
  SZ数据集 & 时空注意力机制-图网络 & 4.05 & 2.74 & 9.21\% \\
  & 时空注意力机制-自适应图网络 & 4.04 & 2.71 & 9.21\% \\
  LA数据集 & 时空注意力机制-图网络 & 6.94 & 3.57 & 10.22\% \\
  & 时空注意力机制-自适应图网络 & 6.34 & 3.03 & 9.06\% \\
  \end{tblr}
\end{table}

另外, 在此部分中, 还将对训练时间进行记录, 研究在二者模型误差较小的情况下, 模型的计算速度, 以便投入实际运用中节省计算资源。主要记录了 100 个 epoch 中每训练 1000batch 数据时, 训练消耗的总时间。

在深圳数据集中, 加入自适应模块共需学习参数量为 174702, 原始模型参数量为 154110。原始模型与改进模型的精度相差不大, 在训练时间上, 原始模型的训练速度比改进后模型的训练速度低 34.76\%以上。在 LA 数据集中, 加入自适应模块共需学习参数量为 289703, 原始模型参数量为 262379。但模型的训练时间, 加入自适应模块训练耗时 \({5.84}{h}\) , 原始模型耗时 \({9.81}{h}\) , 计算速度提高了 \({40.47}\%\)。LA 数据集在 MAE 指标也得到 15.12\%下降。

\begin{table}[htbp]
  \caption{自适应图网络与原模型训练时间对比}
  \centering
  \begin{tblr}{
    hline{1, Z} = {1pt}, 
    hline{2}, 
    width = \textwidth, 
    colspec = {X[1, c]X[1, c]X[1.5, c]X[1.5, c]}, 
    cell{2}{1} = {r=5}{c, m}, 
    cell{7}{1} = {r=5}{c, m}, 
    row{2-Z} = {rowsep = 0pt}
  }
  && {时空注意力机制-\\图网络} & {时空注意力机制-\\自适应图网络} \\
  SZ数据集 & SZ1000steps & 35.40s & 30.58s \\
  & SZ2000steps & 69.57s & 60.13s \\
  & SZ3000steps & 108.56s & 91.21s \\
  & SZ4000steps & 210.90s & 137.89s \\
  &... &... &... \\
  LA数据集 & LA1000steps & 461.93s & 413.49s \\
  & LA2000steps & 1021.12s & 846.13s \\
  & LA3000steps & 1497.50s & 1196.48s \\
  &... &... &... \\
  & LA\_All\_steps & 9.81h & 5.84h \\
  \end{tblr}
\end{table}
